% ============================================================================
% Fichier  : partie3/P03_C13_mots_de_passe.tex
% Titre    : Mots de Passe : Attaques et Défenses
% Auteur   : MINKA MI NGUIDJOÏ Thierry Emmanuel
% Version  : 1.0 -- novembre 2026
% ============================================================================
% Sources : NIST SP 800-63B ; MITRE ATT&CK (T1110 Brute Force et ses
% sous-techniques). Audit des mots de passe de l'Opération LOGONE : cas
% fictif, chiffres cohérents entre eux.
% ============================================================================

\chapter{Mots de Passe : Attaques et Défenses}
\label{chap:mdp}
\chaptermark{Mots de passe}

\epigraph{%
    « Un mot de passe que l'on retient facilement\\
    se devine facilement. »%
}{Principe du cours}

\niveauChapitre{Intermédiaire}{Chapitre~\ref{chap:ad}~; cours SEC3266,
chapitre sur les identités.}

\begin{boxobjectifs}
\begin{itemize}
    \item Distinguer les attaques en ligne et hors ligne.
    \item Expliquer la force brute, le dictionnaire, la pulvérisation et le
          bourrage d'identifiants.
    \item Comprendre pourquoi le hachage et sa lenteur comptent.
    \item Conduire un audit de mots de passe autorisé et en rendre compte
          sans exposer personne.
    \item Proposer une politique conforme au NIST SP~800-63B.
\end{itemize}
\end{boxobjectifs}

\bigskip

% ============================================================================
\section*{Question directrice : combien de mots de passe de LiZbiZ un attaquant retrouverait-il en un jour~?}
% ============================================================================

Après T-01, M.~ETEME demande un audit des mots de passe de l'annuaire, avec
un accord écrit et sous contrôle de Mme~NKENG. NYONG SÉCURITÉ obtient une
copie des empreintes et tente, pendant 24~heures, de retrouver les mots de
passe par des listes de mots courants et des règles de transformation. Le
résultat est un chiffre, et une politique à revoir.

\medskip
\textbf{Plan du chapitre.}
\begin{enumerate}
    \item En ligne ou hors ligne (section~\ref{sec:c13-ligne}).
    \item Les attaques (section~\ref{sec:c13-attaques}).
    \item Le hachage (section~\ref{sec:c13-hachage}).
    \item L'audit de mots de passe (section~\ref{sec:c13-audit}).
    \item La politique (section~\ref{sec:c13-politique}).
\end{enumerate}

\separateur

% ============================================================================
\section{En Ligne ou Hors Ligne}
\label{sec:c13-ligne}
% ============================================================================

\begin{tableaucours}{L{2.6cm}L{4.8cm}L{4.8cm}}{%
    \tableauHeader{Attaque} &
    \tableauHeader{Principe} &
    \tableauHeader{Limite}
}
En ligne & Essais contre le service (portail, messagerie) & Lente,
visible, freinée par le verrouillage et le second facteur \\
Hors ligne & Essais contre une empreinte volée (base, ticket Kerberos) &
Limitée seulement par la puissance de calcul et la lenteur du hachage \\
\end{tableaucours}
\vspace{-0.8em}
\begin{center}{\small\itshape Tableau~13.1. Deux familles d'attaques}\end{center}

% ============================================================================
\section{Les Attaques}
\label{sec:c13-attaques}
% ============================================================================

\begin{tableaucours}{L{3.4cm}L{8.8cm}}{%
    \tableauHeader{Attaque (ATT\&CK T1110)} &
    \tableauHeader{Principe}
}
Force brute & Toutes les combinaisons jusqu'à une longueur donnée \\
Dictionnaire et règles & Mots courants, noms, années, transformés
(majuscule, chiffre, symbole) \\
Pulvérisation & Un ou deux mots de passe courants essayés sur de nombreux
comptes, pour éviter le verrouillage \\
Bourrage d'identifiants & Couples identifiant et mot de passe volés
ailleurs, réutilisés \\
\end{tableaucours}
\vspace{-0.8em}
\begin{center}{\small\itshape Tableau~13.2. Les attaques sur les mots de passe~\cite{mitreattack}}\end{center}

\begin{boxalert}
\textbf{Violation typique.} Croire qu'une règle de complexité (majuscule,
chiffre, symbole) rend un mot de passe fort.

\smallskip
\textbf{Règle~:} « Lizbiz2026! » respecte la règle et figure dans tout
dictionnaire bien conçu. La longueur et l'imprévisibilité comptent~; les
règles de transformation des attaquants connaissent les habitudes humaines.
\end{boxalert}

% ============================================================================
\section{Le Hachage}
\label{sec:c13-hachage}
% ============================================================================

Un système stocke une \emph{empreinte} du mot de passe, pas le mot de passe.
La vitesse à laquelle l'attaquant peut essayer dépend de la fonction de
hachage~: une fonction rapide (comme l'empreinte NT de Windows) permet des
milliards d'essais par seconde sur du matériel courant~; une fonction lente
et salée, conçue pour les mots de passe (bcrypt, scrypt, Argon2), en permet
des milliers. Le sel empêche de casser tous les comptes à la fois.

% ============================================================================
\section{L'Audit de Mots de Passe}
\label{sec:c13-audit}
% ============================================================================

\begin{boxprocedure}
\textbf{Conduire un audit de mots de passe}
\begin{enumerate}
    \item Accord écrit de la direction~; information du délégué du personnel
          selon les règles internes.
    \item Extraction des empreintes par l'administrateur, sur un poste isolé~;
          copie chiffrée.
    \item Cassage hors ligne, durée et méthodes fixées à l'avance.
    \item Résultats agrégés~; aucun mot de passe en clair dans le rapport~;
          liste nominative remise au seul responsable, pour réinitialisation.
    \item Destruction des empreintes et des résultats à la fin.
\end{enumerate}
\end{boxprocedure}

\begin{boxscenario}{Opération LOGONE. Vingt-quatre heures de cassage}
\textbf{Constatations.} 160~comptes actifs après la revue de mars. En
24~heures~: 38~mots de passe retrouvés (23,75~\%), dont 31 dans la
première heure. 12~comptes partagent un même mot de passe avec au moins un
autre compte. 9~mots de passe retrouvés contiennent « lizbiz ». Aucun compte
d'administration n'est retrouvé, sauf celui de T-01, déjà connu.

\medskip
\textbf{Analyse} (confiance élevée). Un attaquant qui volerait les
empreintes ouvrirait près d'un compte sur quatre en une journée. Les
mots de passe retrouvés respectent tous la règle de complexité en vigueur.
\end{boxscenario}

\begin{boxobjet}{A-01 · Mots de passe prévisibles malgré la règle de complexité}
\textbf{Mission complémentaire~:} ce constat relève de l'audit des mots de
passe, rapporté à part des quatorze constats du test (Chapitre~\ref{chap:rapport}).
\textbf{Description~:} 38~mots de passe sur 160 retrouvés en 24~heures.
\textbf{Preuve~:} statistiques agrégées~; liste nominative remise à
M.~KOUNOU. \textbf{Gravité~:} moyenne (élevée en cas de vol des empreintes).
\textbf{Recommandation~:} longueur minimale de 12 à 15~caractères, liste de
mots interdits (nom de l'entreprise, années), contrôle contre les listes
divulguées, second facteur, réinitialisation des 38~comptes.
\textbf{Détection~:} échecs de connexion répartis sur de nombreux comptes
(pulvérisation).
\end{boxobjet}

% ============================================================================
\section{La Politique}
\label{sec:c13-politique}
% ============================================================================

\begin{tableaucours}{L{5.0cm}L{7.2cm}}{%
    \tableauHeader{Recommandation (NIST SP~800-63B)} &
    \tableauHeader{Application chez LiZbiZ}
}
Longueur plutôt que complexité & 12~caractères au moins, phrases de passe
acceptées \\
Comparer aux listes de mots de passe compromis & Contrôle à chaque
changement \\
Pas de changement périodique imposé sans motif & Changement en cas de
soupçon de compromission \\
Limiter les essais & Verrouillage progressif, surveillance \\
Second facteur & Obligatoire pour les accès distants et les comptes à
privilèges \\
\end{tableaucours}
\vspace{-0.8em}
\begin{center}{\small\itshape Tableau~13.3. Une politique conforme au NIST SP~800-63B~\cite{nist80063b}}\end{center}

\begin{boxmetier}{RSSI}
\textbf{Rendre compte sans désigner.} Le RSSI présente l'audit en chiffres
agrégés au comité~; les personnes concernées reçoivent une demande de
changement, accompagnée de conseils, et non un reproche.
\end{boxmetier}

\separateur

\begin{boxresume}
\textbf{Ce qu'il faut retenir de ce chapitre~:}
\begin{itemize}
    \item \textbf{En ligne}~: lent et visible~; \textbf{hors ligne}~: limité
          par le calcul et le hachage.
    \item \textbf{Force brute, dictionnaire, pulvérisation, bourrage}~: quatre
          attaques.
    \item \textbf{Un hachage lent et salé} freine le cassage.
    \item \textbf{Un audit de mots de passe} est autorisé, agrégé, détruit
          à la fin.
    \item \textbf{Longueur, listes de compromis, second facteur}~: la
          politique du NIST.
\end{itemize}
\end{boxresume}

\begin{boxinfo}
\textbf{Points de contrôle de l'élève.} Je sais~: distinguer attaques en
ligne et hors ligne~; nommer les quatre attaques~; expliquer le rôle du
hachage et du sel~; conduire un audit de mots de passe~; rédiger une
politique conforme.
\end{boxinfo}

\bigskip

\begin{boxexo}[title={\ding{45}\quad Exercice~13.1. Le taux \niveaubase}]
Calculez le taux de mots de passe retrouvés en une heure et en 24~heures.
\end{boxexo}

\begin{boxexo}[title={\ding{45}\quad Exercice~13.2. Quelle attaque~? \niveaubase}]
Nommez l'attaque~: « Ete2026! » essayé sur 160~comptes~; les identifiants
d'un site piraté essayés sur la messagerie de LiZbiZ.
\end{boxexo}

\begin{boxexo}[title={\ding{45}\quad Exercice~13.3. Le sel \niveaubase}]
À quoi sert le sel dans le hachage d'un mot de passe~?
\end{boxexo}

\begin{boxexo}[title={\ding{45}\quad Exercice~13.4. Complexité \niveaumoyen}]
Pourquoi « Lizbiz2026! » respecte-t-il la règle de complexité et se
casse-t-il vite~?
\end{boxexo}

\begin{boxexo}[title={\ding{45}\quad Exercice~13.5. Pulvérisation \niveaumoyen}]
Pourquoi la pulvérisation échappe-t-elle au verrouillage de compte~?
Comment la détecter~?
\end{boxexo}

\begin{boxexo}[title={\ding{45}\quad Exercice~13.6. Rapport \niveaumoyen}]
Rédigez la section du rapport sur l'audit de mots de passe, sans exposer
personne.
\end{boxexo}

\begin{boxexo}[title={\ding{45}\quad Exercice~13.7. Réutilisation \niveaumoyen}]
12~comptes partagent un mot de passe. Quel risque~? Que faire~?
\end{boxexo}

\begin{boxexo}[title={\ding{45}\quad Exercice~13.8. Hachage \niveaumoyen}]
Comparez l'empreinte NT et Argon2 du point de vue d'un attaquant hors
ligne.
\end{boxexo}

\begin{boxexo}[title={\ding{45}\quad Exercice~13.9. L'accord \niveauexpert}]
Rédigez l'accord écrit d'audit des mots de passe~: objet, méthodes, durée,
données, destruction, information des salariés.
\end{boxexo}

\begin{boxexo}[title={\ding{45}\quad Exercice~13.10. Politique \niveauexpert}]
Rédigez la nouvelle politique de mots de passe de LiZbiZ et le plan de sa
mise en œuvre.
\end{boxexo}

\bigskip

\begin{boxinfo}
\textbf{Pour aller plus loin.}
\begin{itemize}
    \item NIST SP~800-63B~\cite{nist80063b}~; MITRE ATT\&CK~\cite{mitreattack}.
\end{itemize}
\end{boxinfo}

\begin{boxinfo}
\textbf{Transition.} La Partie~IV passe des systèmes aux applications~: le
Chapitre~\ref{chap:web} apprend à tester une application web.
\end{boxinfo}
